\chapter{Related Work}
\label{chap:related-work}

This chapter compares the inactivity-risk workflow with prior operationalizations
of inactivity, project-health metrics, and security-practice indicators.

\section{OSS Inactivity and Abandonment Operationalizations}
\label{sec:related-inactivity}

Robinson et al. apply survival-style analysis to open-source Python projects and
use revision activity as an observable signal for project survival
\cite{robinson2022survival}. Their results support reasoning about future project
state from activity. A fitted hazard model can yield per-subject predictions, but
their study reports population-level survival behavior. It does not evaluate
fixed-horizon repository predictions against held-out future outcomes as done
here.

Miller et al. find that \(15\%\) of widely used npm packages became abandoned
within six years and that only \(18\%\) of exposed downstream projects removed
the dependency \cite{miller2025abandonment}.\kiref{Claude Opus 5/10} Their
interview study finds that developers often lack clear response strategies
\cite{miller2023wingingit}. Both studies characterize abandonment
retrospectively at ecosystem level. They do not predict a 12-month outcome for an
individual repository from point-in-time features. Their findings motivate an
early-warning triage signal.

Coelho et al.\ train a classifier that identifies unmaintained GitHub projects
from activity features and validate its output with project developers
\cite{coelho2018identifying}. Their follow-up tracks \(2{,}927\) active projects
for one year and proposes a machine-learning metric for the \emph{level} of
maintenance activity. It reports that \(16\%\) became unmaintained
\cite{coelho2020maintained}. At component level, OSSARA aggregates maintenance
predictions to assess the abandonment risk of embedded open-source dependencies
for continuous monitoring \cite{li2022ossara}. These studies show that
maintenance status can be predicted and are the closest predictive comparisons
for this thesis.

The present work differs in its operationalization:

\begin{itemize}
    \item a fixed 12-month horizon evaluated from point-in-time observation snapshots reconstructed from GH~Archive with an explicit temporal leakage boundary;
    \item a multi-signal inactivity label combining human commit, contributor, release, and merged-pull-request activity instead of a single revision-activity threshold or continuous activity index;
    \item probability \emph{calibration} and comparison against strong trivial recency baselines in addition to discrimination metrics;
    \item a separate analysis of the repositories still active at the observation time, where early warning is most decision-relevant;
    \item integration of the calibrated score into a per-repository dependency-triage demonstrator that keeps inactivity prediction and security posture explicitly separate.
\end{itemize}

Issue activity and responsiveness signals are kept as predictors only and are not part of the label.

\section{Repository Health Metrics and Dashboards}
\label{sec:related-health-metrics}

The CHAOSS project provides standardized metrics models, including the Starter
Project Health Metrics Model \cite{chaossStarterProjectHealth}. Goggins et al.
argue that transparent project health requires structured measurement and
interpretation instead of isolated indicators \cite{goggins2021transparent}.
This vocabulary informs the feature engineering.

Dashboards and metrics models primarily describe indicators for human
interpretation. They do not predict a future outcome. Yehudi et al. show that
individual context-free indicators are insufficient for identifying
sustainability across heterogeneous projects \cite{yehudi2023contextfree}. The
model therefore combines observation-time signals and evaluates the combination
empirically. Its output remains a limited operational proxy.

\section{Security-Practice Indicators}
\label{sec:related-security-practice}

OpenSSF Scorecard provides automated security-related checks for open-source
repositories \cite{zahan2023scorecard}. The NIST SSDF and SLSA explain why this
security posture matters for dependencies
\cite{souppaya2022ssdf,slsa2025spec}. The checks describe present security
practices, not future maintenance activity.

OSS Risk Radar therefore presents OpenSSF Scorecard as a parallel
security-posture signal and excludes it from the inactivity-model feature vector
(Section~\ref{sec:method-research-design}). No claim is made that these checks
improve inactivity prediction; they provide complementary triage context.

\section{Responsiveness Measurement Caveats}
\label{sec:related-responsiveness}

The responsiveness features cover issue first-response time and pull-request
response and merge latency. Bot-generated responses can distort these measures
\cite{hasan2023firstresponse}, so they are treated as reconstructed GH Archive
proxies, not canonical CHAOSS metrics
(Section~\ref{sec:background-metric-pitfalls}).

\section{Research Gap}
\label{sec:related-summary}

Prior work operationalizes inactivity and abandonment, standardizes health
metrics, provides security-practice indicators, and predicts maintenance status
\cite{coelho2018identifying,coelho2020maintained,li2022ossara}. The open question
is whether the operationalization in Section~\ref{sec:related-inactivity}
supports a reproducible, forward-looking, \emph{calibrated} per-repository score
within a dependency-triage artifact.
